টুরিং যন্ত্রগুলি তাদের নির্দেশ-সেট দিয়ে নির্ধারিত হয়। এগুলি
পঞ্চকের সসীম সেট (প্রতিটি দশা ও পঠিত প্রতীকের জন্য নতুন দশা,
লিখতে হবে এমন প্রতীক, এবং হেডের চলন নির্দিষ্ট করে)। দশা, প্রতীক ও
চলনকে সংখ্যা দিয়ে প্রকাশ করলে এই পঞ্চকগুলির সসীম সেটের পরিবার
গণনীয়, তাই প্রতিটি টুরিং যন্ত্রের
নির্দেশ-সেটের সঙ্গে একটি সংখ্যা যুক্ত করার উপায় আছে। এমন একটি
নির্দিষ্ট ছক অনুযায়ী টুরিং যন্ত্রের নির্দেশ-সেটের সঙ্গে যে সংখ্যাটি
যুক্ত হয়, তাকে যন্ত্রটির \emph{সূচক} বলে। এইভাবে টুরিং যন্ত্রের
সূচকের কথা বলেই যন্ত্রগুলি নিয়ে পরোক্ষভাবে ``কথা বলতে'' পারি।

টুরিং যন্ত্রের আচরণ নিয়ে একটি গুরুত্বপূর্ণ প্রশ্ন হলো, তারা
শেষ পর্যন্ত থামে কি না। $h(e, n)$ এমন একটি অপেক্ষক ধরি যার মান
$= 1$ হয় যদি সূচক~$e$-বিশিষ্ট টুরিং যন্ত্রটি নিবেশ~$n$ দিয়ে
চালু করলে থামে, এবং অন্যথায় $=0$। একে \emph{থামা-অপেক্ষক} বলে।
থামা-অপেক্ষকটি নিজে টুরিং-গণনসাধ্য কি না, এই প্রশ্নকে বলে
\emph{থামার সমস্যা}। উত্তর হলো না—থামার সমস্যা অসমাধানযোগ্য।
একটি কর্ণ যুক্তির সাহায্যে এটি প্রতিষ্ঠা করা হয়।

``টুরিং যন্ত্র দিয়ে $X$ করা যায় কি'' ধরনের সমস্যার একটি বড়
শ্রেণির মাত্র একটি উদাহরণ হলো থামার সমস্যা। যুক্তিবিদ্যার আরেকটি
কেন্দ্রীয় সমস্যা হলো \emph{প্রথম-ক্রমের যুক্তিবিদ্যার সিদ্ধান্ত
সমস্যা}—এমন কোনো টুরিং যন্ত্র কি আছে যা একটি প্রদত্ত বাক্য বৈধ কি
না তা নির্ণয় করতে পারে? এই বিখ্যাত সমস্যারও নেতিবাচক মীমাংসা
হয়েছে। সিদ্ধান্ত সমস্যা অসমাধানযোগ্য। একটি হ্রাসকরণ যুক্তি দিয়ে
এটি প্রতিষ্ঠা করা হয়। প্রত্যেক টুরিং যন্ত্র~$M$ এবং নিবেশ~$w$-এর
সঙ্গে একটি প্রথম-ক্রমের বাক্য~$T(M, w) \lif E(M, w)$ যুক্ত করা যায়,
যা বৈধ হয় ঠিক তখনই, যখন~$M$ নিবেশ~$w$ দিয়ে চালু করলে থামে।
সিদ্ধান্ত সমস্যা সমাধানযোগ্য হলে এইভাবে সেটি ব্যবহার করে থামার
সমস্যার সমাধান করা যেত।
